\cvsection{Exp\'eriences professionnelles}

\begin{cventries}
  \cventry
  {Doctorat - Universit\'e d'Aix-Marseille - \'Ecole Doctorale Science du vivant \cvwebsite{https://ecole-doctorale-658.univ-amu.fr}{(ED658)}}
  {Laboratoire Dynamics and nanoenvironment of biological membranes \cvwebsite{https://sites.google.com/view/fm4b-lab/home}{(DyNaMo)}}
  {Marseille, France}
  {juin 2022 - mai 2026}
  {
    \begin{cvitems}
      \item {D\'evelopp\'e et automatis\'e des workflows d'acquisition HS-AFM, incluant un logiciel de contr\^ole FPGA personnalis\'e et des pipelines d'analyse post-traitement en Python pour les mesures de forces biologiques}
      \item {Recrut\'e dans le cadre des projets financ\'es par l'ERC \cvwebsite{https://cordis.europa.eu/article/id/456256-the-forces-that-power-immune-defences}{(MechaDynA)}\cvwebsite{https://mission-europe-recherche.fr/fr/projets-laureats/conseil-europeen-recherche-erc/openfmlab}{(OpenFMLab)}}
    \end{cvitems}
    \vspace{1em}
  }

  \cventry
  {Ing\'enieur TEST \& AUTOMATISATION DE LABORATOIRE}
  {Laboratoire Adhesion et Inflammation \cvwebsite{https://labadhesioninflammation.org}{(LAI)}}
  {Marseille, France}
  {mars 2019 - juin 2022}
  {
    \begin{cvitems}
      \item{D\'evelopp\'e des logiciels LabVIEW et NI FPGA pour l'acquisition de donn\'ees \`a haute vitesse en microscopie \`a force atomique rapide (HSAFM)}
      \item {Impl\'ement\'e des algorithmes de contr\^ole pour HSAFM longue port\'ee afin d'am\'eliorer la densit\'e de donn\'ees et le temps total d'acquisition, sous LabVIEW FPGA}
      \item {Exp\'erience pratique en \`electronique: amplificateurs transimp\'edance pour photodiodes, \'electronique de puissance}
      \item {Projet financ\'e par l'ERC \cvwebsite{https://www.europeandissemination.eu/wp-content/uploads/2024/03/MechaDynA-compressed.pdf}{(MechaDynA)}}
    \end{cvitems}
    \vspace{1em}
  }

  \cventry
  {Ing\'enieur d\'eveloppement test}
  {Integrated Micro-electronics, Inc. \cvwebsite{https://beta.global-imi.com}{(IMI)}}
  {Bi\~nan Laguna, Philippines}
  {nov. 2016 - janv. 2019}
  {
    \begin{cvitems}
      \item{D\'evelopp\'e des logiciels et bancs de test automatis\'es sous LabVIEW pour la production, r\'edig\'e la documentation des syst\`emes, accompagn\'e le d\'eploiement \`a l'international, diagnostiqu\'e les probl\`emes de test et optimis\'e les processus de production.}
    \end{cvitems}
    \vspace{1em}
  }

  \cventry
  {Stagiaire ing\'enieur \`electrique}
  {Shell Philippines Exploration B.V (SPEX)}
  {Manille, Philippines}
  {janv. - sept. 2016}
  {
    \begin{cvitems}
      \item{R\'ealis\'e une \`etude de faisabilit\'e sur une centrale CCGT (Combined Cycle Gas Turbine) \`a Malampaya Offshore Gas pour analyser les options d'utilisation locale ou d'injection d'\'electricit\'e sur le r\'eseau national}
    \end{cvitems}
    \vspace{1em}
  }
\end{cventries}
